Data agents answer questions by selecting sources, writing SQL or Python
programs, executing them, and reporting computed values. Auditing these
answers requires connecting the stated answer to the final program, declared
evidence, and data actually read. We present \ledger, a deterministic
post-execution audit layer that canonicalizes declared sources, reconciles
observed reads, replays the final program within a citation-restricted evidence
view, and returns a replay certificate, a localized failure, or an unknown
status. When replay is unavailable, bounded closure over program-reached
operands provides a closure certificate. \ledger also exposes a task-specific
semantic-validation interface, keeping program intent separate from execution
provenance.

We evaluate these capabilities across three increasingly realistic evidence
regimes. On WikiSQL, instantiating the semantic interface with normalized SQL
structure detects all 2,715 controlled semantic perturbations, including 466
that preserve the original execution result, with no false positives on 500
unmodified reference queries. On ConvFinQA and MultiHiertt, \ledger recovers
89.6--93.3\% of reference derivations while rejecting 96.1--97.2\% of
cross-example unsupported answers; standard anomaly detectors reject fewer
than 20\% under the same high-recall calibration. On human-labelled
KramaBench episodes, replay-based selection raises the proportion of delivered
answers judged supported from 28.2\% to 52.4\% while retaining every supported
case; on 148 scored live episodes, it increases the mean official score from
0.307 to 0.429 at 64.2\% coverage. \ledger is deterministic, makes zero model
calls, and adds sub-millisecond bookkeeping on saved traces. These results
show that execution-provenance auditing and task-specific semantic validation
provide complementary evidence for selective delivery and escalation.
